% GENERATED FILE. Edit canonical lessons, then run npm run generate:books.
% canonical-source-hash: fnv1a64:18ddf602848bd355
% canonical-lessons: AR-C228-shahada, AR-C228-istamaa, AR-C228-rasama, AR-C228-ghanna, AR-C228-sawwara

\chapter{To watch, To listen, To draw, To sing, To photograph}
\label{ch:ar-to-watch-to-listen-to-draw-to-sing-to-photograph}
\hlchaptermodality{228}

\begin{chapteropening}
\textbf{By the end of this chapter:} \emph{I can say to watch, to listen, to draw, to sing and to photograph.}

Complete the last lesson of chapter 228: I can say to watch, to listen, to draw, to sing and to photograph.
\end{chapteropening}

\section[\texorpdfstring{\ar{شاهد} (shāhada)}{شاهد (shāhada)}]{\ar{شاهد} (shāhada) — to watch}

\label{lesson:AR-C228-shahada}



\begin{quote}
Before the new one: say the Arabic for to quarrel, then the Arabic for to make up.
\end{quote}

\subsection*{You'll want to know: \ar{شاهد}}
\textbf{\ar{شاهد}} — \emph{shāhada} — \textquotedblleft{}to watch\textquotedblright{}.

\begin{cousinweb}
\textbf{\ar{شاهد}} (\emph{shāhada}) is built on the native Arabic root \emph{sh-h-d}, \textquotedblleft{}to witness, to be present\textquotedblright{} (Lane). \emph{Shāhada} is Form III, \textquotedblleft{}to see with one's own eyes\textquotedblright{}.
\end{cousinweb}

\begin{grammarlens}[title={What you've built}]
One more word for this chapter.
\end{grammarlens}

\subsection*{Your turn}
\begin{itemize}
\raggedright
  \item \emph{Say it:} \emph{shāhada}
  \item \emph{Say it:} \emph{shāhada}, once more
  \item \emph{Say it:} say \emph{shāhada}
  \item \emph{Recall:} say the Arabic for to quarrel, then the Arabic for to make up, then say \emph{shāhada} again
\end{itemize}

\begin{tcolorbox}[breakable,colback=teal!4,colframe=teal!35!black,title={Before you move on}]
What does \ar{شاهد} mean? (\textquotedblleft{}To watch\textquotedblright{}.) Say it once more.
\end{tcolorbox}
\section[\texorpdfstring{\ar{استمع} (istamaʿa)}{استمع (istamaʿa)}]{\ar{استمع} (istamaʿa) — to listen}

\label{lesson:AR-C228-istamaa}



\begin{quote}
Before the new one: say the Arabic for to make up, then the Arabic for to watch.
\end{quote}

\subsection*{You'll want to know: \ar{استمع}}
\textbf{\ar{استمع}} — \emph{istamaʿa} — \textquotedblleft{}to listen\textquotedblright{}.

\begin{cousinweb}
\textbf{\ar{استمع}} (\emph{istamaʿa}) is built on the native Arabic root \emph{s-m-ʿ}, \textquotedblleft{}to hear\textquotedblright{} (Lane). \emph{Istamaʿa} is Form VIII, \textquotedblleft{}to listen\textquotedblright{}.
\end{cousinweb}

\begin{grammarlens}[title={What you've built}]
One more word for this chapter.
\end{grammarlens}

\subsection*{Your turn}
\begin{itemize}
\raggedright
  \item \emph{Say it:} \emph{istamaʿa}
  \item \emph{Say it:} \emph{istamaʿa}, once more
  \item \emph{Say it:} say \emph{istamaʿa}
  \item \emph{Recall:} say the Arabic for to make up, then the Arabic for to watch, then say \emph{istamaʿa} again
\end{itemize}

\begin{tcolorbox}[breakable,colback=teal!4,colframe=teal!35!black,title={Before you move on}]
What does \ar{استمع} mean? (\textquotedblleft{}To listen\textquotedblright{}.) Say it once more.
\end{tcolorbox}
\section[\texorpdfstring{\ar{رسم} (rasama)}{رسم (rasama)}]{\ar{رسم} (rasama) — to draw, to paint}

\label{lesson:AR-C228-rasama}



\begin{quote}
Before the new one: say the Arabic for to watch, then the Arabic for to listen.
\end{quote}

\subsection*{You'll want to know: \ar{رسم}}
\textbf{\ar{رسم}} — \emph{rasama} — \textquotedblleft{}to draw, to paint\textquotedblright{}.

\begin{cousinweb}
\textbf{\ar{رسم}} (\emph{rasama}) is built on the native Arabic root \emph{r-s-m}, \textquotedblleft{}a trace, a mark\textquotedblright{} (Lane). \emph{Rasama} is to trace, to draw.
\end{cousinweb}

\begin{grammarlens}[title={What you've built}]
One more word for this chapter.
\end{grammarlens}

\subsection*{Your turn}
\begin{itemize}
\raggedright
  \item \emph{Say it:} \emph{rasama}
  \item \emph{Say it:} \emph{rasama}, once more
  \item \emph{Say it:} say \emph{rasama}
  \item \emph{Recall:} say the Arabic for to watch, then the Arabic for to listen, then say \emph{rasama} again
\end{itemize}

\begin{tcolorbox}[breakable,colback=teal!4,colframe=teal!35!black,title={Before you move on}]
What does \ar{رسم} mean? (\textquotedblleft{}To draw, to paint\textquotedblright{}.) Say it once more.
\end{tcolorbox}
\section[\texorpdfstring{\ar{غنى} (ghannā)}{غنى (ghannā)}]{\ar{غنى} (ghannā) — to sing}

\label{lesson:AR-C228-ghanna}



\begin{quote}
Before the new one: say the Arabic for to listen, then the Arabic for to draw.
\end{quote}

\subsection*{You'll want to know: \ar{غنى}}
\textbf{\ar{غنى}} — \emph{ghannā} — \textquotedblleft{}to sing\textquotedblright{}.

\begin{cousinweb}
\textbf{\ar{غنى}} (\emph{ghannā}) is built on the native Arabic root \emph{gh-n-y}, \textquotedblleft{}to sing\textquotedblright{} (Lane). \emph{Ghannā} is Form II, \textquotedblleft{}to sing\textquotedblright{}.
\end{cousinweb}

\begin{grammarlens}[title={What you've built}]
One more word for this chapter.
\end{grammarlens}

\subsection*{Your turn}
\begin{itemize}
\raggedright
  \item \emph{Say it:} \emph{ghannā}
  \item \emph{Say it:} \emph{ghannā}, once more
  \item \emph{Say it:} say \emph{ghannā}
  \item \emph{Recall:} say the Arabic for to listen, then the Arabic for to draw, then say \emph{ghannā} again
\end{itemize}

\begin{tcolorbox}[breakable,colback=teal!4,colframe=teal!35!black,title={Before you move on}]
What does \ar{غنى} mean? (\textquotedblleft{}To sing\textquotedblright{}.) Say it once more.
\end{tcolorbox}
\section[\texorpdfstring{\ar{صور} (ṣawwara)}{صور (ṣawwara)}]{\ar{صور} (ṣawwara) — to photograph, to film}

\label{lesson:AR-C228-sawwara}



\begin{quote}
Before the new one: say the Arabic for to draw, then the Arabic for to sing.
\end{quote}

\subsection*{You'll want to know: \ar{صور}}
\textbf{\ar{صور}} — \emph{ṣawwara} — \textquotedblleft{}to photograph, to film\textquotedblright{}.

\begin{cousinweb}
\textbf{\ar{صور}} (\emph{ṣawwara}) is built on the native Arabic root \emph{ṣ-w-r}, \textquotedblleft{}to form, to fashion\textquotedblright{} (Lane). \emph{Ṣawwara} is Form II, \textquotedblleft{}to make an image\textquotedblright{}.
\end{cousinweb}

\begin{grammarlens}[title={What you've built}]
One more word for this chapter.
\end{grammarlens}

\subsection*{Your turn}
\begin{itemize}
\raggedright
  \item \emph{Say it:} \emph{ṣawwara}
  \item \emph{Say it:} \emph{ṣawwara}, once more
  \item \emph{Say it:} say \emph{ṣawwara}
  \item \emph{Recall:} say the Arabic for to draw, then the Arabic for to sing, then say \emph{ṣawwara} again
\end{itemize}

\begin{tcolorbox}[breakable,colback=teal!4,colframe=teal!35!black,title={Before you move on}]
What does \ar{صور} mean? (\textquotedblleft{}To photograph, to film\textquotedblright{}.) Say it once more.
\end{tcolorbox}
